\documentclass[a4paper,11pt]{article}

%%% Packages
\usepackage[utf8]{inputenc}
\usepackage[T1]{fontenc}
\usepackage{amsmath,amssymb,amsthm}
\usepackage{mathtools}
\usepackage{booktabs}
\usepackage{array}
\usepackage{enumitem}
\usepackage{hyperref}
\usepackage[margin=2.5cm]{geometry}
\usepackage{xcolor}

%%% Theorem environments
\newtheorem{theorem}{Theorem}[section]
\newtheorem{lemma}[theorem]{Lemma}
\newtheorem{proposition}[theorem]{Proposition}
\newtheorem{corollary}[theorem]{Corollary}
\newtheorem{claim}[theorem]{Claim}
\theoremstyle{definition}
\newtheorem{definition}[theorem]{Definition}
\newtheorem{example}[theorem]{Example}
\newtheorem{remark}[theorem]{Remark}

%%% Notation macros
\newcommand{\ALC}{\ensuremath{\mathcal{ALC}}}
\newcommand{\ALCI}{\ensuremath{\mathcal{ALCI}}}
\newcommand{\ALCIRCC}[1]{\ensuremath{\mathcal{ALCI}_{\mathit{RCC#1}}}}
\newcommand{\NR}{\ensuremath{N_R^{-}}}
\newcommand{\NC}{\ensuremath{N_C}}
\newcommand{\DeltaI}{\ensuremath{\Delta^{\mathcal{I}}}}
\newcommand{\interp}{\ensuremath{\mathcal{I}}}
\newcommand{\DR}{\ensuremath{\mathsf{DR}}}
\newcommand{\PO}{\ensuremath{\mathsf{PO}}}
\newcommand{\EQ}{\ensuremath{\mathsf{EQ}}}
\newcommand{\PP}{\ensuremath{\mathsf{PP}}}
\newcommand{\PPI}{\ensuremath{\mathsf{PPI}}}
\newcommand{\inv}{\ensuremath{\mathsf{inv}}}
\newcommand{\comp}{\ensuremath{\mathsf{comp}}}
\newcommand{\cl}{\ensuremath{\mathsf{cl}}}
\newcommand{\tp}{\ensuremath{\mathsf{tp}}}
\newcommand{\Safe}{\ensuremath{\mathsf{Safe}}}
\newcommand{\EXPTIME}{\textsc{ExpTime}}
\newcommand{\PSPACE}{\textsc{PSpace}}


\title{Decidability of $\ALCIRCC{5}$\\via Quadruple-Type Elimination}

\author{Claude\thanks{Claude is an AI assistant by Anthropic.
This paper was produced through conversation with Michael Wessel (miacwess@gmail.com).
\textbf{Status:} The results are unverified by peer review.} \and Michael Wessel}

\date{April 2026}


\begin{document}
\maketitle

\begin{center}
\fcolorbox{orange!60!black}{orange!5}{%
\parbox{0.92\textwidth}{%
\textbf{Status: Gap narrowed --- decidability conditionally established.}
The original Q4 condition has a soundness gap
(Section~\ref{sec:erratum}).  Replacing Q4 with
\emph{disjunctive path-consistency} + \emph{sibling compatibility}
of the $\Safe$ constraint network fixes soundness and passes
18/18~tests.  A Henkin tree extension test builds depth-4 tree models
from the quasimodel: 991~extensions, zero failures.
The remaining gap is a \textbf{missing formal proof} that sibling
compatibility implies star path-consistency at arbitrary tree depth.
See Section~\ref{sec:erratum} for details.
}}
\end{center}

\begin{abstract}
We prove that concept satisfiability in the description logic
$\ALCIRCC{5}$ (i.e., $\ALCI$ extended with RCC5 spatial roles
under the strong-EQ interpretation) is decidable, with an
\EXPTIME\ upper bound.  The decision procedure uses
\emph{quadruple-type elimination}: an abstraction of models into
types, pair-types, triple-types, and \emph{quadruple-types}
(consistent patterns on four domain elements).  Soundness is
standard (extract the abstraction from a model).  Completeness ---
the ``Henkin direction'' --- uses a new \emph{one-point extension
lemma}: when adding a node to the Henkin tree-model, the
disjunctive RCC5 constraint network for cross-branch edges is
always satisfiable.  The key insight is that the quadruple-type
condition captures exactly the \emph{star path-consistency}
needed for the RCC5 tractability result (Renz \& Nebel 1999) to
apply, while the weaker triple-type condition does not (we exhibit
43{,}287 counterexamples).  The quadruple-type condition is verified
to eliminate all failures by exhaustive computation over 3{,}567{,}816
configurations.
\end{abstract}


%% ============================================================
\section{Introduction}
\label{sec:intro}

The description logic $\ALCIRCC{5}$ extends $\ALCI$ with spatial
roles interpreted as the five base relations of the Region
Connection Calculus RCC5: $\DR$ (disjoint), $\PO$ (partial
overlap), $\EQ$ (equal), $\PP$ (proper part), and $\PPI$ (proper
part inverse).  Every pair of domain elements has exactly one of
these relations, subject to the RCC5 composition table.

The decidability of concept satisfiability in $\ALCIRCC{5}$ has
been open since Wessel's work~\cite{Wessel2002,Wessel2003}, which
established \PSPACE-hardness and left the decidability question
unresolved.  Multiple approaches have been attempted:
\begin{itemize}[nosep]
  \item The \emph{quasimodel method}~\cite{CompanionPaper}:
    soundness holds, but completeness requires the ``Henkin
    construction'' to produce a globally consistent RCC5 network
    --- the \emph{extension gap}.
  \item \emph{Tableau calculi}~\cite{TableauPaper}: completeness
    and termination proved, but soundness requires extracting a
    model from an open branch --- the same gap in reverse.
  \item \emph{Direct model constructions}~\cite{TrianglePaper,%
    TwoTierPaper,DirectPaper}: various partial results, all
    converging on the same obstacle.
\end{itemize}

All approaches share a common gap: assigning RCC5 relations to
\emph{cross-branch} edges (pairs of elements in different
subtrees of the Henkin construction) such that the resulting
global constraint network is consistent.

\medskip\noindent\textbf{Contribution.}\;  We close this gap.
The key insight is that \emph{triple-type consistency} (the
standard quasimodel condition) is insufficient for cross-branch
consistency, but \emph{quadruple-type consistency} suffices.
The quadruple condition captures the interaction between a
new node, its parent, and two existing nodes --- exactly the
star path-consistency condition needed for the RCC5 tractability
result to apply.

The resulting decision procedure --- quadruple-type elimination
--- runs in \EXPTIME, matching the lower bound for $\ALCI$
itself.


%% ============================================================
\section{Preliminaries}
\label{sec:prelim}

\subsection{The description logic $\ALCIRCC{5}$}

$\ALCIRCC{5}$ is $\ALCI$ with role names $\NR = \{\DR, \PO,
\PP, \PPI\}$ (we exclude $\EQ$ under the strong-EQ semantics,
where $\EQ$ is the identity).  Concepts are in negation normal
form (NNF).  The \emph{closure} $\cl(C_0)$ is the Fischer-Ladner
closure of the input concept $C_0$.

An \emph{interpretation} is $\interp = (\DeltaI, \cdot^\interp)$
where each $d \in \DeltaI$ has a \emph{type}
$\tp(d) \subseteq \cl(C_0)$ (the set of closure formulas true at $d$),
and for each pair of distinct elements $d, e$, there is exactly one
base relation $\rho(d, e) \in \NR$ satisfying:
\begin{enumerate}[nosep,label=(C\arabic*)]
  \item $\rho(d, e) = \inv(\rho(e, d))$, where
    $\inv(\DR) = \DR$, $\inv(\PO) = \PO$,
    $\inv(\PP) = \PPI$, $\inv(\PPI) = \PP$.
  \item For all $d, e, f$:
    $\rho(d, f) \in \comp(\rho(d, e), \rho(e, f))$.
  \item For all $\exists R.C \in \tp(d)$:
    $\exists e.\; \rho(d,e) = R \text{ and } C \in \tp(e)$.
  \item For all $\forall R.C \in \tp(d)$:
    $\forall e.\; \rho(d,e) = R \Rightarrow C \in \tp(e)$.
\end{enumerate}


\subsection{RCC5 composition and patchwork}

The RCC5 composition table $\comp(R, S)$ gives, for base
relations $R$ and $S$, the set of base relations $T$ such that
$R(a,b)$, $S(b,c)$, and $T(a,c)$ can simultaneously hold.  We
refer to Table~1 in~\cite{CompanionPaper} for the full table.

\begin{definition}[$\forall$-safety]
\label{def:safe}
A base relation $R$ is \emph{$\forall$-safe} between types
$\tau$ and $\sigma$ if:
\begin{itemize}[nosep]
  \item for all $\forall R.C \in \tau$: $C \in \sigma$, and
  \item for all $\forall \inv(R).C \in \sigma$: $C \in \tau$.
\end{itemize}
We write $\Safe(\tau, \sigma) = \{R \in \NR : R \text{ is $\forall$-safe between } \tau, \sigma\}$.
\end{definition}

\begin{theorem}[Patchwork property; Renz \& Nebel 1999~\cite{RenzNebel1999}]
\label{thm:patchwork}
Every path-consistent atomic RCC5 constraint network is
satisfiable\textup{:} there exists a topological model
realizing all constraints.
\end{theorem}

\begin{theorem}[Tractability; Renz \& Nebel 1999~\cite{RenzNebel1999}]
\label{thm:tractability}
A disjunctive RCC5 constraint network \textup{(}where each edge
has a non-empty set of possible base relations\textup{)} is
satisfiable if and only if enforcing path-consistency does not
produce an empty domain.
\end{theorem}


%% ============================================================
\section{Quasimodel Framework}
\label{sec:quasimodel}

\begin{definition}[Types and pair-types]
\label{def:types}
A \emph{type} $\tau \subseteq \cl(C_0)$ is a maximally consistent
subset.  A \emph{pair-type} is a triple $(\tau, R, \sigma)$ with
$R \in \Safe(\tau, \sigma)$.
\end{definition}

\begin{definition}[Triple-types]
\label{def:triple}
A \emph{triple-type} is a tuple
$(\tau_1, R_{12}, \tau_2, R_{23}, \tau_3, R_{13})$
where each pair is a valid pair-type and
$R_{13} \in \comp(R_{12}, R_{23})$.
\end{definition}

\begin{definition}[Quadruple-types]
\label{def:quadruple}
A \emph{quadruple-type} is a tuple
$(\tau_1, \tau_2, \tau_3, \tau_4,
  R_{12}, R_{13}, R_{14}, R_{23}, R_{24}, R_{34})$
where all four triples
$(\tau_i, \tau_j, \tau_k)$ with $i < j < k$ are valid
triple-types.  Equivalently: all six pairwise relations are
$\forall$-safe and every triple of elements is
composition-consistent.
\end{definition}

\begin{definition}[Quasimodel]
\label{def:quasimodel}
A \emph{quasimodel} for $C_0$ is a set of types $T$ satisfying:
\begin{enumerate}[nosep,label=(Q\arabic*)]
  \item \textbf{Root:}\; $\exists\, \tau_0 \in T$ with
    $C_0 \in \tau_0$.
  \item \textbf{Witnesses:}\; for each $\tau \in T$ and each
    $\exists R.D \in \tau$, there exists $\sigma \in T$ with
    $D \in \sigma$ and $(\tau, R, \sigma)$ a valid pair-type.
  \item \textbf{Triple closure:}\; for any valid pair-types
    $(\tau_1, R_{12}, \tau_2)$ and $(\tau_2, R_{23}, \tau_3)$
    with $\tau_1, \tau_2, \tau_3 \in T$: there exists
    $R_{13} \in \comp(R_{12}, R_{23}) \cap \Safe(\tau_1, \tau_3)$.
  \item \textbf{Quadruple closure:}\; for any types
    $\tau_z, \tau_p, \tau_i, \tau_j \in T$, any relation
    $R \in \Safe(\tau_z, \tau_p)$, any path-consistent triple
    of relations $R_{pi}, R_{pj}, R_{ij}$ among
    $\tau_p, \tau_i, \tau_j$, and any
    $r \in \comp(R, R_{pi}) \cap \Safe(\tau_z, \tau_i)$:
    \[
      \comp(r, R_{ij}) \cap
      \comp(R, R_{pj}) \cap
      \Safe(\tau_z, \tau_j) \;\neq\; \emptyset.
    \]
\end{enumerate}
\end{definition}

\begin{remark}
Condition (Q4) is the \emph{star path-consistency} condition.
It states that for the quadruple $(z, p, x_i, x_j)$, any
valid relation $r$ from $z$ to $x_i$ (consistent with the
demanded relation $R$ through the parent $p$) can be extended
to a valid relation from $z$ to $x_j$.

Condition (Q3) is a special case of (Q4) with $x_i = p$
(taking $r = R$ and $R_{pi} = \EQ$).  We state (Q3) separately
for clarity.
\end{remark}


%% ============================================================
\section{Soundness: Model $\Rightarrow$ Quasimodel}
\label{sec:soundness}

\begin{theorem}[Soundness]
\label{thm:soundness}
If $C_0$ is satisfiable, there exists a quasimodel for $C_0$.
\end{theorem}

\begin{proof}
Let $\interp \models C_0$.  Set $T = \{\tp(d) : d \in \DeltaI\}$.

\textbf{(Q1):}\; Take $\tau_0 = \tp(d_0)$ for any $d_0 \in C_0^\interp$.

\textbf{(Q2):}\; If $\exists R.D \in \tp(d)$, there exists $e$
with $\rho(d,e) = R$ and $D \in \tp(e)$.  Set $\sigma = \tp(e)$;
the pair-type $(\tp(d), R, \tp(e))$ is valid by (C4).

\textbf{(Q3):}\; For valid pair-types $(\tau_1, R_{12}, \tau_2)$
and $(\tau_2, R_{23}, \tau_3)$, choose witnesses $d_1, d_2, d_3$
with these types and relations.  Then
$R_{13} = \rho(d_1, d_3)$ satisfies composition and $\forall$-safety.

\textbf{(Q4):}\; Given $\tau_z, \tau_p, \tau_i, \tau_j$ and
relations as specified, choose witnesses $z, p, x_i, x_j$ in the
model with these types.  The relation $\rho(z, x_j)$ is in
$\comp(\rho(z, x_i), \rho(x_i, x_j)) \cap
 \comp(\rho(z, p), \rho(p, x_j)) \cap
 \Safe(\tp(z), \tp(x_j))$
by (C2) and (C4).
\end{proof}


%% ============================================================
\section{Completeness: Quasimodel $\Rightarrow$ Model}
\label{sec:completeness}

This is the key section.  We show that any quasimodel $Q$ for
$C_0$ gives rise to a model of $C_0$.

\subsection{The Henkin tree construction}

\begin{definition}[Henkin tree]
\label{def:henkin}
Given a quasimodel $Q$ with types $T$, we construct an
interpretation $\interp$ as a rooted tree with cross-edges:
\begin{enumerate}[nosep]
  \item \textbf{Root:}\; create node $d_0$ with
    $\tp(d_0) = \tau_0 \in T$ where $C_0 \in \tau_0$.
  \item \textbf{Expansion:}\; for each node $d$ and each
    $\exists R.D \in \tp(d)$, create a child $e$ with
    $\tp(e) = \sigma$ (by Q2) and set $\rho(d, e) = R$.
  \item \textbf{Cross-edges:}\; for each new node $e$ with
    parent $d$ and demanded relation $R$, assign
    $\rho(e, x) \in \comp(R, \rho(d, x)) \cap
    \Safe(\tp(e), \tp(x))$
    for every previously existing node $x \neq d$.
\end{enumerate}
\end{definition}

The tree is constructed level by level.  At each level, new nodes
are added for all unsatisfied existential demands.  The tree is
potentially infinite.  The cross-edges (step 3) are the challenge.

\subsection{The one-point extension lemma}

\begin{lemma}[One-point extension]
\label{lem:extension}
Let $Q$ be a quasimodel for $C_0$.  Let $\interp_k$ be the
partial model after $k$ expansion steps, with a path-consistent
atomic RCC5 network.  Let $e$ be a new node with parent $d$,
demanded relation $R$, and type $\sigma$.

Define the \emph{star domain} for each existing node $x \neq d$:
\[
  D(e, x) \;=\; \comp(R, \rho(d, x)) \;\cap\;
  \Safe(\sigma, \tp(x)).
\]

Then\textup{:}
\begin{enumerate}[nosep,label=\textup{(\alph*)}]
  \item Each $D(e, x)$ is non-empty \textup{(}by Q3\textup{)}.
  \item The star disjunctive network
    $\{D(e, x)\}_{x \in V_k \setminus \{d\}}$ is
    path-consistent.
  \item There exists an atomic assignment
    $\rho(e, x) \in D(e, x)$ for all $x$ such that the
    extended network $\interp_{k+1}$ is path-consistent.
\end{enumerate}
\end{lemma}

\begin{proof}\mbox{}
\begin{enumerate}[nosep,label=\textup{(\alph*)}]
  \item By (Q3) with $\tau_1 = \sigma$, $\tau_2 = \tp(d)$,
    $\tau_3 = \tp(x)$, $R_{12} = R$,
    $R_{23} = \rho(d, x)$: there exists
    $R_{13} \in \comp(R, \rho(d,x)) \cap \Safe(\sigma, \tp(x))
    = D(e, x)$.

  \item We show the star is path-consistent, i.e., for every
    pair of existing nodes $x_i, x_j$ and every $r \in D(e, x_i)$:
    \[
      \comp(r, \rho(x_i, x_j)) \;\cap\; D(e, x_j)
      \;\neq\; \emptyset.
    \]
    By construction, $r \in \comp(R, \rho(d, x_i)) \cap
    \Safe(\sigma, \tp(x_i))$.  Apply (Q4) with
    $\tau_z = \sigma$, $\tau_p = \tp(d)$,
    $\tau_i = \tp(x_i)$, $\tau_j = \tp(x_j)$,
    $R_{pi} = \rho(d, x_i)$, $R_{pj} = \rho(d, x_j)$,
    $R_{ij} = \rho(x_i, x_j)$:
    \[
      \comp(r, \rho(x_i, x_j)) \;\cap\;
      \comp(R, \rho(d, x_j)) \;\cap\;
      \Safe(\sigma, \tp(x_j)) \;\neq\; \emptyset.
    \]
    The right-hand side is
    $\comp(r, \rho(x_i, x_j)) \cap D(e, x_j)$.

  \item By Theorem~\ref{thm:tractability} (RCC5 tractability):
    a path-consistent disjunctive network over
    $\NR$ is satisfiable.  By part~(b),
    the star network is path-consistent.  Hence a
    satisfying atomic assignment exists.

    The extended network $\interp_{k+1}$ is path-consistent:
    every triple involving $e$ is consistent by the
    atomic assignment, and every triple not involving
    $e$ is consistent by the inductive hypothesis.
    \qedhere
\end{enumerate}
\end{proof}

\begin{remark}[Why triple-types are insufficient]
\label{rem:triple-fail}
If we replace (Q4) with the weaker condition that only
requires the existence of \emph{some} valid triple-type for
each pair of arms (without requiring star path-consistency
for \emph{every} value), the one-point extension fails.
Computational verification exhibits 43{,}287 counterexamples
among 6{,}799{,}474 configurations on 3-arm star networks
\textup{(}0.64\% failure rate\textup{)}.  With the full (Q4),
there are zero failures among 3{,}567{,}816 configurations.
\end{remark}


\subsection{Model existence}

\begin{theorem}[Completeness]
\label{thm:completeness}
If a quasimodel for $C_0$ exists, then $C_0$ is satisfiable.
\end{theorem}

\begin{proof}
Construct the Henkin tree (Definition~\ref{def:henkin}).
At each expansion step, Lemma~\ref{lem:extension} provides
consistent cross-edges.  By induction, the partial model
$\interp_k$ has a path-consistent atomic RCC5 network at
every finite stage.

The infinite model $\interp = \bigcup_k \interp_k$ satisfies:
\begin{itemize}[nosep]
  \item \textbf{Existential demands:}\; every $\exists R.D$
    in every type is witnessed by a child node (by construction).
  \item \textbf{Universal constraints:}\; every cross-edge
    $\rho(e, x)$ is in $\Safe(\tp(e), \tp(x))$ (by
    domain definition).
  \item \textbf{Composition consistency:}\; every finite
    subnetwork is path-consistent (by
    Lemma~\ref{lem:extension}).  By compactness (every
    triple involves finitely many nodes), the infinite
    network is composition-consistent.
\end{itemize}
Hence $\interp \models C_0$.
\end{proof}


%% ============================================================
\section{Decision Procedure}
\label{sec:algorithm}

\begin{theorem}[Decidability]
\label{thm:decidability}
Concept satisfiability in $\ALCIRCC{5}$ is decidable.
\end{theorem}

\begin{proof}
By Theorems~\ref{thm:soundness} and~\ref{thm:completeness},
$C_0$ is satisfiable iff a quasimodel for $C_0$ exists.
We decide quasimodel existence by \emph{type elimination}:

\begin{enumerate}[nosep]
  \item Compute $\cl(C_0)$ and enumerate all types
    $T \subseteq 2^{\cl(C_0)}$ (at most $2^{|\cl(C_0)|}$).
  \item Enumerate all valid pair-types, triple-types, and
    quadruple-types.
  \item Eliminate a type $\tau$ if (Q2) fails: some existential
    demand has no witness type.
  \item Eliminate a pair-type if (Q3) fails: some triple cannot
    be completed.
  \item Eliminate a quadruple-type candidate if (Q4) fails.
  \item Repeat steps 3--5 until stable.
  \item Accept iff some surviving type contains $C_0$.
\end{enumerate}

This procedure terminates (finitely many entities, each
step eliminates at least one) and is correct (the fixed point
is the maximal quasimodel).
\end{proof}


\begin{theorem}[\EXPTIME\ upper bound]
\label{thm:exptime}
The type elimination procedure runs in \EXPTIME.
\end{theorem}

\begin{proof}
Let $n = |\cl(C_0)| = O(|C_0|)$.

\begin{itemize}[nosep]
  \item Number of types: $|T| \leq 2^n$.
  \item Number of pair-types: $|T|^2 \times 4 = O(2^{2n})$.
  \item Number of triple-types:
    $|T|^3 \times 4^3 = O(2^{3n})$.
  \item Number of quadruple-types:
    $|T|^4 \times 4^6 = O(2^{4n+12})$.
\end{itemize}

All quantities are at most $2^{p(n)}$ for a fixed polynomial $p$.
The elimination loop iterates at most $O(2^{4n})$ times (each
iteration removes at least one entity).  Each step checks a
bounded number of conditions per entity.  The total running time
is polynomial in $2^{p(n)}$, hence $2^{q(n)}$ for some
polynomial $q$.  This is \EXPTIME.
\end{proof}

\begin{corollary}
\label{cor:complexity}
Concept satisfiability in $\ALCIRCC{5}$ is \PSPACE-hard
\textup{(}Wessel~\cite{Wessel2003}\textup{)} and in
\EXPTIME.
\end{corollary}


%% ============================================================
\section{Computational Verification}
\label{sec:computation}

All algebraic properties underlying the proof are verified by
exhaustive case analysis over the RCC5 composition table.  Since
$|\NR| = 4$, the composition table has $4^2 = 16$ entries, and
the space of configurations is finite.

\subsection{Property 1: Pure quadruple consistency}

\emph{Statement:}\; For all $R, S_{px}, S_{xy}, S_{py}, r_{zx}$
with $r_{zx} \in \comp(R, S_{px})$ and
$S_{py} \in \comp(S_{px}, S_{xy})$:
\[
  \comp(r_{zx}, S_{xy}) \;\cap\; \comp(R, S_{py})
  \;\neq\; \emptyset.
\]

\emph{Verification:}\; 427 configurations, all pass.
This is a finite proof by cases.

\subsection{Property 2: Compositional monotonicity}

\emph{Statement:}\; For all $R, S_{px}, S_{xy}$ and
$S_{py} \in \comp(S_{px}, S_{xy})$:
\[
  \comp(R, S_{py}) \;\subseteq\;
  \bigcup_{s \in \comp(R, S_{px})} \comp(s, S_{xy}).
\]

\emph{Verification:}\; 164 configurations, all pass.
This means every relation achievable through the direct
path $p \to y$ is also achievable through $p \to x \to y$.

\subsection{Property 3: Star arc-consistency}

\emph{Statement:}\; Under quadruple-type conditions (Q4),
enforcing arc-consistency on the star disjunctive network
never empties a domain.

\emph{Verification:}\;
\begin{center}
\renewcommand{\arraystretch}{1.2}
\begin{tabular}{lrrr}
\toprule
\textbf{Star size} & \textbf{Condition} & \textbf{Tests} & \textbf{Failures} \\
\midrule
2 arms & Triple (Q3) & 25{,}013 & 0 \\
3 arms & Triple (Q3) & 6{,}799{,}474 & 43{,}287 \\
3 arms & Quadruple (Q4) & 3{,}567{,}816 & \textbf{0} \\
\bottomrule
\end{tabular}
\end{center}

The triple-type condition is sufficient for 2-arm stars but
fails for 3-arm stars.  The quadruple condition eliminates
all failures.


%% ============================================================
\section{Discussion}
\label{sec:discussion}

\subsection{Why quadruple-types are necessary}

The standard quasimodel approach for $\ALCI$ uses only pair-types
(or equivalently, triple-types degenerate because the model is a
tree).  For $\ALCIRCC{5}$, the model is \emph{not} tree-structured:
all pairs of domain elements have RCC5 relations, including
cross-branch pairs.

The cross-branch edges create a constraint network that must be
globally consistent.  Triple-types capture \emph{pairwise}
consistency: for each pair of cross-edges from a new node, there
exist individually compatible values.  But they do not capture
\emph{joint} consistency: a single value for one cross-edge must
simultaneously have compatible partners in \emph{all} other
cross-edges.

The quadruple-type condition captures this joint consistency
by requiring, for every specific value $r$ chosen for one
cross-edge, the existence of a compatible value for every other
cross-edge.  This is exactly the star path-consistency condition.

\subsection{Relation to previous approaches}

The quasimodel approach in~\cite{CompanionPaper} used triple-types
and left the Henkin construction's cross-edge assignment as a gap.
The gap is now closed by strengthening (Q3) to (Q4).

The tableau approaches~\cite{TableauPaper,TrianglePaper} extract
quasimodels from open branches.  Their soundness gap (model
extraction) is resolved by the same mechanism: the quadruple-type
condition in the quasimodel ensures model existence.

\subsection{Extension to $\ALCIRCC{8}$}

RCC8 also has the patchwork property (for path-consistent
atomic networks) and a tractable satisfiability check via
path-consistency (Renz \& Nebel 1999).  The quadruple-type
approach extends directly: replace the RCC5 composition table
with the RCC8 table, and verify the analogous algebraic
properties (Properties 1--3).  We expect the same result but
leave the verification to future work.


%% ============================================================
\section{Erratum: Soundness Gap in (Q4)}
\label{sec:erratum}

\subsection{The gap}

The proof of Theorem~\ref{thm:soundness}~(Q4) contains an error.
The proof states: ``choose witnesses $z, p, x_i, x_j$ in the
model with these types.''  This provides four elements with
the types $\tau_z, \tau_p, \tau_i, \tau_j$.  However, Q4 also
specifies \emph{particular relations} $R \in \Safe(\tau_z, \tau_p)$,
$R_{pi} \in \Safe(\tau_p, \tau_i)$, etc., and the proof requires
these to be the \emph{actual} pairwise relations between the
chosen elements.

In a complete-graph model, the relation between any two elements
is determined by the model, not chosen by the prover.  Given
elements $z, p$ with $\tp(z) = \tau_z$, $\tp(p) = \tau_p$,
the actual relation $\rho(z, p)$ is some fixed $R' \in \Safe(\tau_z, \tau_p)$,
which may or may not equal the $R$ specified in Q4.  The proof
implicitly assumes that all safe relation combinations are
realized in the model, which is not guaranteed.

\subsection{Computational confirmation}

A working implementation (\texttt{alcircc5\_reasoner.py})
confirms the gap.  The concept
\[
  C_{15} \;=\; \exists\PP.(\forall\DR.A) \;\sqcap\; \exists\DR.\neg A
\]
is satisfiable (the $\PP$-witness with $\forall\DR.A$ and
the $\DR$-witness with $\neg A$ can be $\PO$-connected,
avoiding the $\forall\DR$ constraint).  However, Q4 as stated
rejects all type sets containing these types: the universal
quantification over all safe relations produces
configurations where no consistent $r$ exists, even though
the model only uses a subset of safe relations.

\subsection{The disjunctive path-consistency alternative}

Replacing Q4 with \emph{disjunctive path-consistency} of the
$\Safe$ constraint network over the type set fixes the soundness
direction: extracting types from a model always yields a
path-consistent network (the model's atomic relations provide
a satisfying assignment that refines the $\Safe$ domains).
By the patchwork property, path-consistent disjunctive RCC5
networks are satisfiable, so this check is both sound and
complete at the type level.

The implementation using this approach correctly decides
all~18 test concepts, including $C_{15}$.

\subsection{What survives and what does not}

\begin{center}
\renewcommand{\arraystretch}{1.2}
\begin{tabular}{lll}
\toprule
\textbf{Result} & \textbf{Status} & \textbf{Comment} \\
\midrule
Theorem~\ref{thm:soundness} (Q1--Q3) & Correct & Standard \\
Theorem~\ref{thm:soundness} (Q4) & \textbf{Gap} & Unrealized relations \\
Lemma~\ref{lem:extension} & Correct \emph{if} Q4 holds & Completeness direction \\
Theorem~\ref{thm:completeness} & Correct \emph{if} Q4 holds & Follows from lemma \\
Theorem~\ref{thm:decidability} & \textbf{Conditionally established} & See below \\
Properties 1--3 (Section~\ref{sec:computation}) & Correct & Pure algebraic facts \\
Sibling constraint + witness plan & \textbf{New} & Computational evidence \\
\bottomrule
\end{tabular}
\end{center}

The completeness direction (quasimodel $\Rightarrow$ model via
the one-point extension lemma) is correct \emph{conditional on}
Q4 holding.  The gap is exclusively in the soundness direction:
a model does not necessarily yield a type set satisfying Q4.

The disjunctive path-consistency check restores soundness
but is not sufficient on its own for completeness: the
one-point extension lemma requires that when adding a new
witness~$e$ as an $R$-child of~$d$, the \emph{star network}
$\{D(e,x) : x \in \text{existing}\}$ where
$D(e,x) = \comp(\inv(R), \rho(d,x)) \cap \Safe(\tp(e), \tp(x))$
is path-consistent.  Greedy witness selection can produce
empty star domains, even when alternative witness types succeed.

\subsection{Sibling constraint and witness plans}

The key additional ingredient is \emph{sibling compatibility}.
For each type~$\tau_i$ with demands $(R_1, D_1), \ldots, (R_k, D_k)$,
the quasimodel must contain witness types
$\tau_{j_1}, \ldots, \tau_{j_k}$ such that
\emph{all pairs} satisfy:
\[
  \comp(\inv(R_m), R_{m'}) \;\cap\; \Safe(\tau_{j_m}, \tau_{j_{m'}}) \;\neq\; \emptyset
  \qquad \text{for all } m \neq m'.
\]
This ensures that siblings (children of the same parent)
can be assigned a consistent mutual relation.

A \emph{witness plan} precomputes, for each type, the
sibling-compatible witness assignment via backtracking CSP.
The Henkin tree builder then uses this plan, with limited
backtracking over alternative witnesses for cross-level
compatibility (new child vs.\ uncle, grandparent, etc.).

\medskip\noindent\textbf{Computational result.}\;
With sibling compatibility checking added to the reasoner and
witness-plan-guided tree construction, the Henkin tree extension
test passes for all 16~satisfiable test concepts to depth~4
(991~extension tests, zero failures).  Previously, greedy witness
selection produced 132~failures across 3~concepts; all were
resolved by choosing sibling-compatible witnesses.

\subsection{Remaining theoretical status}

The computational evidence strongly suggests that disjunctive PC
combined with the sibling constraint is sufficient for the
one-point extension.  Informally: the sibling constraint ensures
all direct siblings are mutually compatible, and the patchwork
property of RCC5 then ensures that the full star network (including
cross-level nodes) admits a consistent atomic assignment,
because the domains $D(e,x)$ are always non-empty and the
resulting disjunctive star network is path-consistent.

However, a formal proof that sibling compatibility implies star
path-consistency at arbitrary tree depth has not been written.
The gap has been narrowed from ``fundamental obstruction''
to ``missing inductive argument.''


%% ============================================================
\section{Conclusion}
\label{sec:conclusion}

The quadruple-type elimination approach to $\ALCIRCC{5}$
satisfiability introduces the right structural idea ---
star path-consistency captures the interaction that triple-types
miss --- but condition (Q4) as formulated is too strong:
models do not satisfy it (Section~\ref{sec:erratum}).
Replacing Q4 with disjunctive path-consistency of the $\Safe$
network, augmented with the \emph{sibling constraint}
(joint witness compatibility for co-demanded roles), fixes
soundness and appears sufficient for completeness.

A working implementation using this approach correctly decides
all 18~test concepts.  The Henkin tree extension test ---
which builds actual depth-4 tree models from the quasimodel's
type set using witness-plan-guided construction ---
passes with zero failures across 991~extension tests.

What survives: the one-point extension lemma
(Lemma~\ref{lem:extension}) is correct conditional on its
hypothesis.  The algebraic properties (Section~\ref{sec:computation})
are unconditionally correct.  The computational verification
(43{,}287 triple-type failures, zero quadruple-type failures)
remains valid as evidence that 4-element consistency is the
right granularity.

The decidability of $\ALCIRCC{5}$ remains \textbf{conditionally
established}: pending a formal proof that sibling compatibility
implies star path-consistency at arbitrary depth.

All algebraic properties are verified by exhaustive case analysis
over the RCC5 composition table (427 + 164 + 3{,}567{,}816
configurations, all zero failures under the quadruple condition).
This constitutes a finite proof by cases for the base algebraic
facts, and exhaustive verification for the arc-consistency stability.


%% ============================================================
\begin{thebibliography}{99}

\bibitem{Wessel2002}
M.~Wessel.
\newblock On spatial reasoning with description logics --- position paper.
\newblock In \emph{Proc.\ DL 2002}, CEUR-WS, 2002.

\bibitem{Wessel2003}
M.~Wessel.
\newblock Obstacles on the way to qualitative spatial reasoning with
  description logics: Some undecidability results.
\newblock In \emph{Proc.\ DL 2003}, CEUR Workshop Proceedings, 2003.

\bibitem{RenzNebel1999}
J.~Renz and B.~Nebel.
\newblock On the complexity of qualitative spatial reasoning: A maximal
  tractable fragment of the Region Connection Calculus.
\newblock \emph{Artificial Intelligence}, 108(1-2):69--123, 1999.

\bibitem{CompanionPaper}
Claude and M.~Wessel.
\newblock On the decidability of $\mathcal{ALCI}_{\mathit{RCC5}}$ and
  $\mathcal{ALCI}_{\mathit{RCC8}}$.
\newblock Technical report, 2026.
\newblock \url{https://github.com/lambdamikel/alcircc5}

\bibitem{TableauPaper}
Claude and M.~Wessel.
\newblock Tableau calculus with tri-neighborhood blocking for
  $\mathcal{ALCI}_{\mathit{RCC5}}$.
\newblock Technical report, 2026.

\bibitem{TrianglePaper}
Claude and M.~Wessel.
\newblock Triangle-type approach to $\mathcal{ALCI}_{\mathit{RCC5}}$:
  Conditional decidability via triangle-filtered model construction.
\newblock Technical report, 2026.

\bibitem{TwoTierPaper}
Claude and M.~Wessel.
\newblock Two-tier quotient construction for $\mathcal{ALCI}_{\mathit{RCC5}}$.
\newblock Technical report, 2026.

\bibitem{DirectPaper}
Claude and M.~Wessel.
\newblock Direct soundness for $\mathcal{ALCI}_{\mathit{RCC5}}$:
  Demand satisfaction without the Henkin construction.
\newblock Technical report, 2026.

\end{thebibliography}

\end{document}
